%
% teil3.tex -- Punkte, Linien und Kreise in der gaussschen Ebene
%
\section{Punkte, Linien und Kreise in der gaussschen Ebene
\label{komplex:section:objekte}}
\kopfrechts{Die gaussschen Ebene}
Bevor die Konstruktion gleichseitiger Dreiecke entwickelt werden kann,
ist es nützlich zu sehen,
wie sich elementare geometrische Objekte --- Strecken, Geraden, Kreise ---
direkt in der Sprache komplexer Zahlen ausdrücken lassen.
Die Unterabschnitte führen Strecken, Geraden samt Mittelsenkrechten,
Kreise und schliesslich Kollinearität und Winkelmessung ein.

%
% Strecken als Differenzen komplexer Zahlen
%
\subsection{Strecken als Differenzen komplexer Zahlen
\label{komplex:subsection:strecken}}
Eine Strecke zwischen zwei Punkten $z_1$ und $z_2$ wird in komplexer
Schreibweise vollständig durch die Differenz $z_2 - z_1$ beschrieben.
Diese Differenz enthält gleichzeitig die Länge der Strecke und ihre
Richtung.
Die Länge ist $|z_2 - z_1|$, die Richtung wird durch das Argument
$\arg(z_2 - z_1)$ festgelegt.

Diese kompakte Codierung ist der eigentliche Grund,
warum komplexe Zahlen für geometrische Konstruktionen so geeignet sind.
Eine Strecke wird zu einer einzigen Zahl, mit der man weiterrechnen kann.

%
% Geraden, Mittelsenkrechte und Parameterdarstellung
%
\subsection{Geraden, Mittelsenkrechte und Parameterdarstellung
\label{komplex:subsection:geraden}}
Eine Gerade durch zwei Punkte $z_1$ und $z_2$ lässt sich parametrisch
schreiben als
\begin{equation}
z(t)
=
z_1 + t \, (z_2 - z_1) ,
\qquad
t \in \mathbb{R} .
\label{komplex:eq:gerade}
\end{equation}
Für $t = 0$ erhält man $z_1$,
für $t = 1$ erhält man $z_2$,
und für $t = 1/2$ den Mittelpunkt der Strecke.
Mit $t \in [0, 1]$ beschreibt die Formel die Strecke selbst,
mit $t \in \mathbb{R}$ die ganze Gerade.

Die \emph{Mittelsenkrechte} zweier Punkte $z_1$ und $z_2$ besteht aus allen
\index{Mittelsenkrechte}%
Punkten $z$, die zu beiden den gleichen Abstand haben:
\begin{equation}
|z - z_1|
=
|z - z_2| .
\label{komplex:eq:mittelsenkrechte}
\end{equation}
Quadrieren wir beide Seiten und nutzen die Identität
$|w|^2 = w \cdot \overline{w}$, so wird daraus
\[
(z - z_1)(\overline{z} - \overline{z_1})
=
(z - z_2)(\overline{z} - \overline{z_2}) .
\]
Ausmultiplizieren liefert auf beiden Seiten den Term $z \overline{z}$,
der sich kürzt; die übrigen Terme ordnen wir nach $z$ und $\overline{z}$
und erhalten
\begin{equation}
(\overline{z_2} - \overline{z_1}) \, z
+ (z_2 - z_1) \, \overline{z}
=
|z_2|^2 - |z_1|^2 .
\label{komplex:eq:mittelsenkrechte-linear}
\end{equation}
Das ist eine lineare Gleichung in $z$ und $\overline{z}$, die genau die
Mittelsenkrechte beschreibt.

%
% Kreise als Niveaumengen des Betrags
%
\subsection{Kreise als Niveaumengen des Betrags
\label{komplex:subsection:kreise}}
\index{Niveaumenge}%
Der Kreis mit Mittelpunkt $z_0$ und Radius $r$ besteht aus allen Punkten
$z$ mit
\begin{equation}
|z - z_0|
=
r .
\label{komplex:eq:kreis}
\end{equation}
Diese Gleichung ist eine direkte Übersetzung der geometrischen Definition.
In ihrer Knappheit liegt ihr Wert:
Während der Kreis in kartesischen Koordinaten zwei reelle Variablen $x$
und $y$ benötigt, beschreibt eine einzige komplexe Gleichung die ganze
Punktmenge.

%
% Kollinearitaet und Winkel
%
\subsection{Kollinearität und Winkel
\label{komplex:subsection:winkel}}
Drei Punkte $z_1, z_2, z_3$ sind genau dann \emph{kollinear},
\index{kollinear}%
wenn der Quotient
\begin{equation}
\frac{z_3 - z_1}{z_2 - z_1}
\label{komplex:eq:quotient}
\end{equation}
reell ist.
Allgemein gibt das Argument dieses Quotienten den Winkel an,
unter dem die Strecke $z_1 z_3$ gegenüber der Strecke $z_1 z_2$ gedreht
ist.
Damit verschmelzen Winkelberechnungen und Quotientenbildung zu derselben
Operation.
Diese Beobachtung ist der Schlüssel für die in
Abschnitt~\ref{komplex:section:beweis} geführten Beweise.
